In a two-colouring of \(K_n\), are there at most \(n^2/4\) edges
which belong to no monochromatic triangle?
Call such edges \emph{uncovered}, and let \(h(n)\) be their maximum
possible number. The exact known answer is
\[
 h(n)=
 \begin{cases}
 \binom n2,&n\leq5,\\
 10,&n=6,\\
 \lfloor n^2/4\rfloor,&n\geq7.
 \end{cases}
\]
Thus the literal assertion for every \(n\) has small counterexamples;
the intended large-\(n\) statement is true.
We reconstruct the proof for \(n\geq10\) from Keevash and Sudakov,
\emph{On the number of edges not covered by monochromatic copies of
a fixed graph}, Theorem 1.1 and Proposition 2.1,
\url{https://people.math.ethz.ch/~sudakovb/nim-edges.pdf}.
Their verified finite cases \(n=7,8,9\) are used externally.
See \url{https://www.erdosproblems.com/639}.
